\documentclass[dvipdfmx]{jlreq}
\usepackage{silence}
\WarningFilter{caption}{Unknown document class}
\usepackage{comment}
\usepackage{mathtools, amssymb, amsthm}
\usepackage{pxrubrica}
\usepackage{graphicx}
\usepackage{caption}
\usepackage{qrcode}
\usepackage{pgfplots}
\usepackage{otf}
\usetikzlibrary{intersections}
\pgfplotsset{compat=1.18}

\title{$\pi > 3.14$ の証明過程}
\author{79期 2-B 12番 \CID{13748}清一誓}
\date{}

\begin{document}
\maketitle
\section{はじめに}
今回は、円周率の下からの近似を行っていこうと思います。わたしは前々から円周率を自分で近似したいと思っていたので、かなり頑張って書いていきます。
$\pi > 3$からどんどん強くしていく感じで計算します。お察しの通り円に内接する多角形で近似していきます。それでは、いきます。
\section{本題：円周率の下からの近似}
\subsection{$\pi > 3$ の証明}
まず、$\pi > 3$の証明をしていく。
\begin{proof}
半径1の円に正六角形が内接していることを考える。このとき、この正六角形の周長は円の周長より短いことを利用する。

半径 $r = 1$ の円の周長は $2\pi \times 1 = 2\pi$ である。

一方、半径1の円に内接する正六角形は、中心から各頂点へ直線を引くことで、1辺の長さが1の合同な正三角形6個に分割できる。
したがって、正六角形の周長は$1 \times 6=6$である。

この正六角形の周長は円の周長よりも短いことより、次の不等式が成り立つ。
\[ 2\pi>6 \]

両辺を2で割ると、円周率 $\pi$ に関する次の下界が得られる。
\[ \pi>3 \]

以上より示せた。
\end{proof}
まあこのぐらいはな実は学校の教科書にも載っているくらいなので、次に行きます。先は長そうです。
\subsection{正六角形→正八角形}
頂点の数は増やせば増やすほど近似はもちろん良くなっていくので、次は計算しやすそうな正八角形で行きます。
\begin{proof}
半径1の円に正八角形が内接していることを考える。このとき、この正八角形の周長は円の周長より短いことを利用する。

半径 $r = 1$ の円の周長は $2\pi \times 1 = 2\pi$ である。

一方、この一辺の長さは余弦定理より
\[ \sqrt{1+1-2\cos \frac{\pi}{4}}=\sqrt{2-\sqrt{2}} \]

この正八角形の周長は円の周長よりも短いことより、次の不等式が成り立つ。
\[ 2\pi>8\sqrt{2-\sqrt{2}} \quad (\approx 6.122934918\dots) \]

両辺を 2 で割ると、次のような不等式が成り立つ事がわかる。
\[ \pi>4\sqrt{2-\sqrt{2}}>3.061 \]
以上より示せた。
\end{proof}
これであの有名な東大の一行問題である「円周率が3.05より大きいことを示せ。」には答えられますね。次に行きましょう。

なお、$\pi>3$を示すのには周長を使うと正六角形で証明できるのに対し、面積で近似すると正十二角形を使わないといけないことから、周長で近似したほうが効率が良いと考えられるため、今後も周長で近似していく。
\subsection{三平方の定理の利用}
座標平面上の円周上に格子点がたくさん出てくる円を利用します。今回は半径325にしましょう。（半径5で一度やってみたが全然だめだったので5:12:13と7:24:25の2つを利用することにしました。8:15:17ではなく7:24:25を採用したのは最小公倍数を小さくするためです。）
\begin{proof}
半径$325$の円$x^2+y^2=325^2$を考える。

三平方の定理より、$(3,4,5)$、$(5,12,13)$、$(7,24,25)$の各辺の長さをそれぞれ$65$倍、$25$倍、$13$倍することで、円周上の点
\[
(195,260),\ (125,300),\ (91,312)
\]
が得られる。また、座標を入れ替えた点も同じ円周上にある。

したがって、第1象限において、円周上の点
\[
\begin{aligned}
&(0,325),\ (91,312),\ (125,300),\\
&(195,260),\ (260,195),\ (300,125),\\
&(312,91),\ (325,0)
\end{aligned}
\]
をこの順に結んだ折れ線を考える。

各線分の長さを三平方の定理で求めると、
\[
\begin{aligned}
\sqrt{91^2+13^2}&=65\sqrt{2},\\
\sqrt{34^2+12^2}&=10\sqrt{13},\\
\sqrt{70^2+40^2}&=10\sqrt{65},\\
\sqrt{65^2+65^2}&=65\sqrt{2},\\
\sqrt{40^2+70^2}&=10\sqrt{65},\\
\sqrt{12^2+34^2}&=10\sqrt{13},\\
\sqrt{13^2+91^2}&=65\sqrt{2}
\end{aligned}
\]
となる。

よって、この折れ線の長さは
\[
65\sqrt{2}+10\sqrt{13}+10\sqrt{65}+65\sqrt{2}+10\sqrt{65}+10\sqrt{13}+65\sqrt{2}
=195\sqrt{2}+20\sqrt{13}+20\sqrt{65}
\]
である。

円の対称性より、この折れ線を各象限に配置してできる内接多角形の周長は
\[
4(195\sqrt{2}+20\sqrt{13}+20\sqrt{65})
\]
となる。内接多角形の周長は円周の長さより短いため、
\[
650\pi>4(195\sqrt{2}+20\sqrt{13}+20\sqrt{65})
\]
が成り立つ。

両辺を$650$で割ると、
\[
\begin{aligned}
\pi
&>\frac{2(195\sqrt{2}+20\sqrt{13}+20\sqrt{65})}{325}\\
&=\frac{390\sqrt{2}+40\sqrt{13}+40\sqrt{65}}{325}\
&\approx3.1330943\dots>3.133
\end{aligned}
\]
となる。

以上より、$\pi>3.133$が示せた。
\end{proof}
あぁ！惜しいですね。一回半径を17倍して8:15:17の直角三角形を利用することも考えたのですが、この感じで行くと3.14には届かなさそうだったので、別のアプローチを考えます、というか試します。あまりにも有名な手法だからです。「内接する正多角形の頂点の数をどんどん倍にしていく」です。これなら半角の公式を用いて理論上いくらでも改良できます。それでは始めます。正六角形はもうやったので、正十二角形から始めます。
\subsubsection{正十二角形を用いた証明}
\begin{proof}
半径$1$の円に内接する正十二角形を考え、その一辺の長さを$a_{12}$とする。

半径$1$の円に内接する正六角形の一辺の長さは$1$である。正$n$角形の一辺の長さを$a_n$とすると、半角の公式より
\[
a_{2n}=\sqrt{2-\sqrt{4-a_n^2}}
\]
が成り立つ。

したがって、
\[
a_{12}=\sqrt{2-\sqrt{4-1^2}}
=\sqrt{2-\sqrt{3}}
\]
となる。

正十二角形の周長は$12a_{12}$であり、円周の長さ$2\pi$より短いため、
\[
2\pi>12\sqrt{2-\sqrt{3}}
\]
が成り立つ。

両辺を$2$で割ると、
\[
\pi>6\sqrt{2-\sqrt{3}}
\approx3.105828541\dots
\]
となる。

以上より、正十二角形を用いて円周率の下界を改良できた。
\end{proof}
どんどん行きます。次は正二十四角形です。
\subsubsection{正二十四角形を用いた証明}
\begin{proof}
半径$1$の円に内接する正二十四角形を考え、その一辺の長さを$a_{24}$とする。

正十二角形の一辺の長さは
\[
a_{12}=\sqrt{2-\sqrt{3}}
\]
であった。半角の公式より、
\[
\begin{aligned}
a_{24}
&=\sqrt{2-\sqrt{4-a_{12}^2}}\\
&=\sqrt{2-\sqrt{2+\sqrt{3}}}
\end{aligned}
\]
となる。

正二十四角形の周長は$24a_{24}$であり、円周の長さ$2\pi$より短いため、
\[
2\pi>24\sqrt{2-\sqrt{2+\sqrt{3}}}
\]
が成り立つ。

両辺を$2$で割ると、
\[
\pi>12\sqrt{2-\sqrt{2+\sqrt{3}}}
\approx3.132628613\dots
\]
となる。

以上より、正二十四角形を用いて円周率の下界をさらに改良できた。
\end{proof}
次は正四十八角形です。
\subsubsection{正四十八角形を用いた証明}
\begin{proof}
半径$1$の円に内接する正四十八角形を考え、その一辺の長さを$a_{48}$とする。

正二十四角形の一辺の長さは
\[
a_{24}=\sqrt{2-\sqrt{2+\sqrt{3}}}
\]
であった。半角の公式より、
\[
\begin{aligned}
a_{48}
&=\sqrt{2-\sqrt{4-a_{24}^2}}\\
&=\sqrt{2-\sqrt{2+\sqrt{2+\sqrt{3}}}}
\end{aligned}
\]
となる。

正四十八角形の周長は$48a_{48}$であり、円周の長さ$2\pi$より短いため、
\[
2\pi>48\sqrt{2-\sqrt{2+\sqrt{2+\sqrt{3}}}}
\]
が成り立つ。

両辺を$2$で割ると、
\[
\begin{aligned}
\pi
&>24\sqrt{2-\sqrt{2+\sqrt{2+\sqrt{3}}}}\\
&\approx3.139350203\dots
\end{aligned}
\]
となる。

以上より、正四十八角形を用いて円周率の下界をさらに改良できた。
\end{proof}
惜しいですね。次は正九十六角形です。
\subsubsection{正九十六角形を用いた証明}
\begin{proof}
半径$1$の円に内接する正九十六角形を考え、その一辺の長さを$a_{96}$とする。

正四十八角形の一辺の長さは
\[
a_{48}=\sqrt{2-\sqrt{2+\sqrt{2+\sqrt{3}}}}
\]
であった。半角の公式より、
\[
\begin{aligned}
a_{96}
&=\sqrt{2-\sqrt{4-a_{48}^2}}\\
&=\sqrt{2-\sqrt{2+\sqrt{2+\sqrt{2+\sqrt{3}}}}}
\end{aligned}
\]
となる。

正九十六角形の周長は$96a_{96}$であり、円周の長さ$2\pi$より短いため、
\[
2\pi>96a_{96}
\]
が成り立つ。

両辺を$2$で割ると、
\[
\begin{aligned}
\pi&>48\sqrt{2-\sqrt{2+\sqrt{2+\sqrt{2+\sqrt{3}}}}}\\
&\approx3.14103195089053\dots>3.14
\end{aligned}
\]
となる。

以上より、$\pi>3.14$が示せた。
\end{proof}
来ました...!　$\pi>3.14$が示せました！
\section{気づいたこと}
ところで、$a_{6\cdot2^n}$って規則的な形をしていますよね。
\[
a_{6\cdot2^n}=\sqrt{2-\underbrace{\sqrt{2+\sqrt{2+\dots+\sqrt{2+\sqrt{3}}}}}_{n-1\text{個の}2}}
\]
が成り立っています。これを利用として$\pi$の近似式を作りたくなったのです。やってみましょう。\section{$\pi$の近似式}

ここまで、正十二角形、正二十四角形、正四十八角形、正九十六角形と、頂点の数を2倍ずつ増やすことで円周率の下界を改良してきた。

ここで、先ほど気づいた規則性を利用して、一般に正$6\cdot 2^n$角形を用いたときの近似式を考えてみる。ただし、$n$は1以上の整数とする。

まず、入れ子の平方根を簡潔に表すため、
\[
b_0=\sqrt{3},\qquad
b_{k+1}=\sqrt{2+b_k}\quad(k\geq 0)
\]
と定める。

すると、これまでの計算結果から、半径1の円に内接する正$6\cdot 2^n$角形の一辺の長さ$a_{6\cdot 2^n}$は
\[
a_{6\cdot 2^n}=\sqrt{2-b_{n-1}}
\]
と表せる。

実際、$n=1$のときは
\[
a_{12}=\sqrt{2-\sqrt{3}},
\]
$n=2$のときは
\[
a_{24}=\sqrt{2-\sqrt{2+\sqrt{3}}},
\]
となり、これまで求めた式と一致する。

したがって、正$6\cdot 2^n$角形の周長は
\[
6\cdot 2^n\sqrt{2-b_{n-1}}
\]
である。この多角形は半径1の円に内接しているため、その周長は円周の長さ$2\pi$より短い。

よって、
\[
2\pi>6\cdot 2^n\sqrt{2-b_{n-1}}
\]
となる。両辺を2で割れば、
\[
\boxed{\pi>3\cdot 2^n\sqrt{2-b_{n-1}}}
\]
が得られる。

ここで、
\[
\boxed{p_n=3\cdot 2^n\sqrt{2-b_{n-1}}}
\]
とおくと、$p_n$は円周率$\pi$を下から近似する数列となる。つまり、$n$を指定すれば、入れ子の平方根を計算することで$\pi$の近似値を求められる。

例えば、これまでの計算結果をまとめると、次のようになる。

\[
\begin{array}{c|c|c}
n & \text{正多角形} & p_n\\ \hline
1 & \text{正十二角形} & 3.105828541\ldots\\
2 & \text{正二十四角形} & 3.132628613\ldots\\
3 & \text{正四十八角形} & 3.139350203\ldots\\
4 & \text{正九十六角形} & 3.141031950\ldots
\end{array}
\]

この表からも、頂点の数を増やすほど近似値が$\pi$に近づいていることがわかる。

では、$n$を限りなく大きくしたとき、この数列は本当に$\pi$に収束するのだろうか（もちろん直感的には明らかだが）。

正$N$角形の一辺の長さは、半径1の円を考えると
\[
a_N=2\sin\frac{\pi}{N}
\]
と表せる。したがって、
\[
\begin{aligned}
p_n
&=3\cdot 2^n a_{6\cdot 2^n}\\
&=6\cdot 2^n
\sin\frac{\pi}{6\cdot 2^n}\\
&=\pi\frac{
\sin\left(\dfrac{\pi}{6\cdot 2^n}\right)
}{
\dfrac{\pi}{6\cdot 2^n}
}.
\end{aligned}
\]

ここで、$n\to\infty$とすると、
\[
\frac{\pi}{6\cdot 2^n}\to 0
\]
となる。また、よく知られた極限
\[
\lim_{x\to 0}\frac{\sin x}{x}=1
\]
を用いれば、
\[
\begin{aligned}
\lim_{n\to\infty}p_n
&=\pi\lim_{n\to\infty}
\frac{
\sin\left(\dfrac{\pi}{6\cdot 2^n}\right)
}{
\dfrac{\pi}{6\cdot 2^n}
}\\
&=\pi
\end{aligned}
\]
が得られる。

以上より、
\[
\boxed{
\pi=\lim_{n\to\infty}
3\cdot 2^n\sqrt{2-b_{n-1}}
}
\]
が成り立つ。ただし、$b_0=\sqrt{3}$、
$b_{k+1}=\sqrt{2+b_k}$である。

つまり、入れ子の平方根を用いて表された数列を考えることで、円周率を下から近似し、その極限として$\pi$を表すことができた。

今回の計算では、正九十六角形を用いることで$\pi>3.14$を示せた。しかし、頂点の数をさらに増やしていけば、より精度の高い近似値が得られる。最初は$\pi>3$という簡単な不等式から始まったが、正多角形の周長と半角の公式を利用することで、円周率の近似を段階的に改良できることがわかった。

最後に、かっこよくするために、$b_n$を用いずに書き下した式の形を書こうと思う。
\[
\boxed{
\pi=\lim_{n\to\infty}
3\cdot 2^n
\sqrt{
2-\underbrace{
\sqrt{2+\sqrt{2+\cdots+\sqrt{2+\sqrt{3}}}}
}_{n-1\text{個の }2}
}
}
\]
\section{おわりに}
ついに$\pi$の近似式まで出すことができました。自分も式の形に規則性があるとはじめは予想できず、とても驚きました。なぜだかとても嬉しいです。ちなみに先生は気づきましたでしょうか、冒頭の僕の名前の兼清の兼が旧字体の\CID{13748}になっていることに。このためだけにotfパッケージを利用しました。だんだんLaTeX組版にも慣れてきました。これからもいっぱいレポート書きます！最後に、今回のTeXソースコードのGithubリンクのQRコードを乗せておきます。
\end{document}
